\title{When Users Set the Agent's Priorities but Often Depart From Its
Recommendation: Explainability, Anthropomorphism, and Cognitive Forcing in
Agentic AI Decision Support}

\author{Anonymous Author(s)}

\begin{abstract}
Explainability is commonly treated as a safeguard against overreliance,
because users are expected to inspect the evidence behind a recommendation
and decide accordingly. In agentic decision support the same assistant
assembles the evidence, interprets the case, and recommends an action, so
the explanation is encountered as part of a socially delivered
recommendation, and cognitive forcing has been proposed as a corrective. We
report a $2 \times 2 \times 2$ hiring experiment ($N=211$) varying
provenance-based explainability cues, anthropomorphic delivery cues, and
human intervention checkpoints, separating assigned conditions, post-task
perceptions, and logged behaviour. Advancement under the study's caution
criterion occurred for 159 of the 211 participants, and no assigned cue was
significantly associated with it. Participants who experienced the
assistant as more explainable also experienced it as more anthropomorphic
and more trustworthy, and post-task trust was associated with advancement.
Pre-recommendation document access was available in every condition and was
used by 55.0\% of participants without distinguishing the decisions. In
checkpoint conditions, 70 participants selected a caution-relevant
evaluation priority, which deterministically produced a \emph{Hold}
recommendation reporting the corresponding evidence gap, yet 46 of them
advanced the candidate. Making users' evaluative priorities consequential
for the agent therefore produced only limited alignment between the
resulting recommendation and the final human action.
\end{abstract}

\keywords{Agentic AI, Explainable AI, Anthropomorphism, Trust, Cognitive
Forcing, Human Intervention, Reliance}

\maketitle

\section{Introduction}

Agentic AI systems increasingly assemble evidence, interpret cases,
recommend actions, and adapt their workflow in response to user input, and
this creates a delegation problem for human-agent interaction, because
providing users with evidence and with formal decision authority does not
by itself ensure that either is used effectively. Explainability is
commonly treated as a safeguard against overreliance, since explanations
are expected to help users understand a recommendation, assess whether it
fits the available evidence, and justify a decision in relation to that
evidence \cite{lee_trust_2004, schemmer_appropriate_2023}. Prior work
shows, however, that explanations do not necessarily improve the detection
of flawed advice when the supporting evidence is not examined
\cite{bucinca_trust_2021, vasconcelos_explanations_2023,
poursabzi-sangdeh_manipulating_2021}, and that they rarely support the
verification required for complementary human--AI performance
\cite{fok_search_2024}. In agentic interaction the explanation is delivered
by the same assistant that interprets the case, and socially fluent
delivery can raise perceived trustworthiness beyond what the evidence
supports \cite{araujo_living_2018, kramer_tricking_2025,
stinkeste_linguistic_2025}. Cognitive forcing was proposed as a response to
this problem, since interrupting the sequence from recommendation to
decision is intended to make users engage with the evidence rather than
accept the recommendation on the strength of its delivery
\cite{bucinca_trust_2021}. Human intervention checkpoints operationalise
that logic at the workflow level, since they allow users to specify what
the agent should examine before a recommendation is generated
\cite{horvitz_principles_1999}. Steering an agent, however, is not the same
as inspecting evidence, and neither is the same as deciding accordingly,
which is of concern in delegated workflows because retained decision
authority does not by itself ensure substantive oversight.

\section{Method}

We conducted a controlled $2 \times 2 \times 2$ between-subjects online
experiment varying provenance-based explainability cues (absent or
present), anthropomorphic delivery cues (low or high), and human
intervention checkpoints (HICs; absent or present). Participants were
recruited through Prolific with prescreening for English fluency and
current or recent experience in human resources, talent acquisition, or
personnel management. Qualtrics randomly assigned participants to one of the
eight conditions.\footnote{The study was approved by [committee and reference
withheld for anonymity]. Participants provided informed consent and
received [payment] for approximately [duration] minutes.} The merged
dataset contained 224 rows. Four duplicate records were removed by
retaining the most complete and most recent matched session for each
participant identifier, and records without a passed attention check, a
matched interaction log, or a final decision were then excluded, leaving
211 participants with 24 to 29 per condition.

Participants acted as recruiters evaluating a fictional candidate for a
Strategic Talent Operations Partner role, and all received the same company
context, role summary, screening-policy summary, and candidate CV. Before
the recommendation, participants could open the full role description and
screening policy through controls available in every condition. HIC-present
participants then encountered a pre-recommendation evaluation-priority
checkpoint. The assistant presented its assessment and recommendation,
HIC-present participants could request a targeted review of a selected
issue, and participants then reported how settled their judgement was and
made a final decision from \emph{Reject}, \emph{Advance to human
interview}, or \emph{Hold for further review}. Perceived explainability,
perceived anthropomorphism, and trust were measured afterwards, so that the
questionnaire could not sensitise participants to the focal constructs
before their evidence access and decision behaviour had been recorded.

The assistant retrieved and organised evidence from the fixed documents,
evaluated the candidate against role and policy requirements, generated one
of the permitted recommendations, and adapted its assessment to
participant-selected priorities. Pre-authored recommendation and follow-up
cards held the evidential content and permitted outcomes stable across
participants. Explainability-present conditions added provenance links
connecting assessment claims to inspectable document sections
\cite{qi_model_2024}. High-anthropomorphism conditions used first-person,
cooperative, and socially fluent wording, whereas low-anthropomorphism
conditions used report-style wording \cite{araujo_living_2018,
stinkeste_linguistic_2025}, and because the manipulation also varied
interactional framing in checkpoint prompts and action labels it should be
understood as a bundle of social-delivery cues rather than warmth alone.

At the first checkpoint, HIC-present participants could optionally select
up to three of six evaluation priorities: independent ownership,
stakeholder communication, transferable experience, structured evaluation
or screening experience, operational coordination, and growth potential.
Independent ownership and structured evaluation or screening experience
were designated caution-relevant, because the CV did not provide evidence
satisfying those role-critical requirements, and selecting either
deterministically produced a recommendation of \emph{Hold} rather than
\emph{Advance}, whereas other selections adapted the assessment rationale
without changing the recommendation. This was a programmed response to
self-selected input rather than an independently generated change of view.
At the second checkpoint, participants could request a targeted review
after seeing the recommendation. HIC exposure was therefore a two-stage
bundle, and checkpoint use was analysed as workflow steering rather than as
evidence verification.

The CV showed relevant recruitment-coordination experience but did not
document independent screening ownership, structured application of
evaluation criteria, or responsibility for progression decisions, so
\emph{Hold} was treated as the study-defined caution-consistent response
and \emph{Advance} was coded as advancement under weak support. This coding
represents the study's policy-based criterion rather than an independently
validated standard. Pre-recommendation full-document access recorded
whether a participant opened the complete role description or screening
policy before the recommendation, and citation inspection was recorded
separately and was possible only where provenance cues were present.
Perceived explainability (four items), perceived anthropomorphism (four
items), and trust (six items) were measured after the decision using
seven-point agreement scales ranging from strongly disagree to strongly
agree and were computed as scale means, with alphas of $.86$, $.93$, and
$.91$, respectively [insert verified measure citations]. Assigned-condition
differences were estimated with an additive logistic model containing the
three factors. Additive analysis-of-variance models assessed the manipulation
checks, and ordinary least-squares models estimated the post-task associations
among the 195 participants with complete perception measures. The indirect
association used 10,000 percentile-bootstrap resamples. Analyses not specified
before data collection are exploratory. Full measure wording, participant
characteristics, and full-factorial sensitivity models will be reported in
the supplementary material.

\section{Results}

In additive manipulation-check models, explainability-cue assignment was
associated with higher perceived explainability, $F(1,199)=4.04$,
$p=.046$, partial $\eta^{2}=.020$, and anthropomorphic-cue assignment was
associated with higher perceived anthropomorphism, $F(1,202)=5.57$,
$p=.019$, partial $\eta^{2}=.027$. Among the 195 participants with complete
perception measures, perceived explainability was associated with perceived
anthropomorphism ($b=0.811$, SE $=0.097$, $p<.001$), and both perceived
explainability ($b=0.653$, SE $=0.058$, $p<.001$) and perceived
anthropomorphism ($b=0.197$, SE $=0.037$, $p<.001$) were associated with
trust, with an estimated indirect association of $b=0.160$, bootstrap
SE $=0.044$, 95\% CI $[0.079, 0.250]$. Because these measures were
collected together after the decision, they describe post-interaction
associations rather than a temporally ordered process.

Advancement under the study's caution criterion was the dominant outcome,
since 159 of the 211 participants, 75.4\%, selected \emph{Advance}. In the
additive assigned-condition model, advancement was not significantly
associated with explainability cues, OR $=0.82$, 95\% CI $[0.44, 1.54]$,
$p=.538$, anthropomorphic cues, OR $=0.98$, 95\% CI $[0.52, 1.84]$,
$p=.948$, or checkpoint presence, OR $=0.78$, 95\% CI $[0.42, 1.46]$,
$p=.439$. The intervals do not exclude modest effects in either direction.
Post-task trust was associated with advancement after adjustment for the
three assigned factors, OR $=1.64$, 95\% CI $[1.14, 2.35]$, $p=.008$,
though this association is not interpreted as showing that trust preceded
the decision.

Before the recommendation was presented, 116 participants, 55.0\%, opened
at least one full source document through controls available in every
condition. Of these, 83, 71.6\%, advanced the candidate, compared with 76
of the 95, 80.0\%, who did not, and the association was not significant,
Fisher exact $p=.199$. Opening a document records evidence access rather
than reading depth, so access alone did not reliably distinguish the
decisions.

Of the 104 participants in HIC-present conditions, 91, 87.5\%, selected at
least one evaluation priority, and 70, 67.3\%, selected a caution-relevant
priority, for whom the deterministic policy produced a \emph{Hold}
recommendation in every case. After receiving that recommendation, 21
participants, 30.0\%, chose \emph{Hold}, 3, 4.3\%, chose \emph{Reject}, and
46, 65.7\%, chose \emph{Advance}. Most participants on this self-selected
path therefore advanced the candidate despite receiving a cautious
recommendation generated in response to their own evaluation priorities.
Because entry into the path followed participants' own selections rather
than random assignment, this is a within-path descriptive result and does
not estimate the causal effect of receiving a \emph{Hold} recommendation.

\section{Discussion and Contribution}

The behavioural result identifies a boundary of this implementation of
workflow-level cognitive forcing. Human-selected priorities were
consequential for the agent, since caution-relevant selections changed its
evidence emphasis and deterministically produced a \emph{Hold}
recommendation, but they did not govern the final human decision, because
46 of the 70 participants on that path still advanced the candidate.
Participants selected the priorities themselves and the recommendation
followed a fixed policy rule, so the finding shows that a consequential
opportunity to steer an agent can coexist with subsequent departure from
its recommendation, rather than demonstrating that checkpoints generally
fail to improve decision quality.

The perception-layer findings complement this account. Assigned cues
determined what the interface presented, whereas the post-task measures
captured how those cues were experienced, and participants who experienced
the assistant as more explainable also experienced it as more
anthropomorphic and more trustworthy, with greater trust associated with
advancement. The questionnaire was administered after the task to avoid
priming attention to the focal constructs, so these relationships are
interpreted as post-interaction associations rather than as evidence that
perceived explainability produced trust or that trust caused advancement.

Two design features constrain the interpretation. \emph{Hold} was defined
as the caution-consistent response from researcher-authored task materials
rather than from an independently validated standard, and the
final-decision interface required participants choosing \emph{Hold} to
select an additional unresolved issue whereas \emph{Advance} and
\emph{Reject} required no equivalent justification, which may have
discouraged \emph{Hold} \cite{schemmer_appropriate_2023}. Future studies
should impose equivalent justification requirements across all decision
options, validate the decision criterion independently, collect an unaided
judgement before the recommendation is presented, and manipulate a
reconciliation step when a final decision conflicts with a recommendation
produced from the participant's own priorities.

\bibliographystyle{ACM-Reference-Format}
\bibliography{HAI}

\end{document}
